\chapter{Distributing Inferred Specifications}
\label{ch:distribution}

The two preceding chapters established that specifications can be inferred with
acceptable fidelity and that they carry measurable value to a downstream
consumer. Both rest on an assumption that is rarely satisfied in practice: that
the consumer has the source tree the inferrer ran over. This chapter removes
the assumption. It develops and evaluates a mechanism for carrying inferred JML
contracts across the boundary that separates source code from compiled artefact
--- into Java class files and into REST service descriptions --- so that a
consumer with only a JAR or only an HTTP interface can recover the
specification the inferrer produced. The chapter answers RQ3.

\section{The Distribution Gap}
\label{sec:dist-gap}

Recent work, including the study of Chapter~\ref{ch:testgen}, has shown that
inferred JML specifications materially improve downstream tooling. That work
shares an unstated assumption: the consumer of the specification --- the test
generator, the integrated development environment, the verifier, the language
model reading an API at inference time --- has access to the same source the
inferrer processed. The assumption is rarely true. Java libraries are
distributed as JAR files containing compiled bytecode, not source. REST
endpoints expose only their HTTP signatures. gRPC services expose only their
protocol-buffer descriptors. Source code is the exception in a consumer's world,
not the norm. A specification that lives only in the source tree is a
specification available only to the team that owns the source --- which is
precisely the team least in need of it, since they can read the code.

The conceptual question of whether specifications \emph{should} travel with
compiled artefacts has a long history in the Design by Contract and JML
traditions~\cite{meyer1992design, leavens2006design, jml}. The practical
question this chapter addresses is narrower and, in a sense, more demanding:
can the \emph{standard, unmodified} Java and REST toolchain be made to carry
specifications, so that enrichment requires no change to how artefacts are
built, signed, published, or consumed? A mechanism that required a custom
class loader, a patched compiler, or a bespoke service framework would not be
adopted; a mechanism that rides inside the existing artefact formats might be.
The research questions are accordingly framed around losslessness, overhead,
throughput, and toolchain compatibility.

\begin{description}
  \item[RQ3a (Losslessness).] Can inferred JML contracts be embedded in
        compiled Java artefacts and reproduced byte-for-byte from the embedded
        bytes, with no source dependency?
  \item[RQ3b (Overhead).] What is the bytecode-size cost of the embedding, and
        how does it scale across libraries of different shapes?
  \item[RQ3c (Throughput).] Is embedding and reading fast enough to support
        per-pull-request continuous-integration workflows on production-scale
        libraries?
  \item[RQ3d (Compatibility).] Does an OpenAPI extension family integrate with
        existing OpenAPI tooling without modification?
\end{description}

\subsection{Why Existing Channels Fall Short}

Section~\ref{sec:bg-distribution} reviewed the existing channels; the gap each
leaves is worth restating concretely here. JML specification files (\texttt{.jml}
stubs) are the verifier's native format but travel as separate sidecar files,
absent from the compiled JAR and invisible to any tool that does not look for
them. Runtime assertion checking compiles specifications into executable checks
inside the bytecode, carrying them into the artefact but only in executable
form, restricted to runtime-evaluable properties, and at the cost of altering
the program's behaviour. The documentation field of an OpenAPI description
carries prose for a human, not a structured contract for a tool. None of these
carries a structured, behaviour-preserving, toolchain-transparent specification
inside the standard artefact that any reader can recover without the source.
That is the gap this chapter fills, on both the bytecode side and the service-
interface side.

\section{Format Design}
\label{sec:dist-format}

The format is the contract between the embedder and any future consumer; it
must be precise enough that an alternative embedder or reader could be written
from the specification alone. The bytecode side uses a runtime-retained Java
annotation; the service side uses an OpenAPI extension family; a sidecar JAR
covers the cases neither can.

\subsection{The Annotation Type}

The embedded specification is carried by a single annotation per program
element, retained at runtime and targetable at methods, constructors, fields,
and types:

\begin{lstlisting}
package com.jml.spec;

@Retention(RetentionPolicy.RUNTIME)
@Target({ElementType.METHOD, ElementType.CONSTRUCTOR,
         ElementType.FIELD, ElementType.TYPE})
public @interface JmlSpecs {
    String[] requires()      default {};
    String[] ensures()       default {};
    String[] assignable()    default {};
    String[] loopInvariant() default {};
    /** Each entry is "ExceptionType|condition". */
    String[] signals()       default {};
    String version()         default "2";
}
\end{lstlisting}

A single annotation carries every clause kind as a separate string array. Within
a member, array order is significant for \texttt{requires}: preconditions are
emitted in source order, which is load-bearing for well-definedness chains,
since a precondition asserting an array element's value depends on earlier
preconditions establishing the array's non-nullness and the index's
in-boundsness. Unrecognised members are skipped on read, which preserves forward
compatibility: a future clause kind added as a new member is ignored by an old
reader rather than breaking it.

\subsection{Three Size Optimisations}

A naive encoding --- one annotation per clause, every clause spelled out in full
--- is wasteful, and the format applies three optimisations whose justification
is empirical, grounded in the shape distribution of real inferred specifications.

The first is the \emph{single packed annotation}. An earlier format version used
Java's repeatable-annotation mechanism, emitting one annotation per clause with
per-clause kind, text, order, and version members wrapped in a container. This
is the maximally verbose shape, paying the per-clause annotation framing --- a
type reference, an element-value-pair structure, an end marker --- for every
clause. The current format collapses this to one annotation per element with the
clause kinds as array members, eliminating the per-clause framing entirely.

The second is \emph{default omission for the frame condition}. Every inferred
specification in the corpus carries \texttt{assignable \textbackslash nothing},
the engine's default for methods that appear pure. The format omits the
\texttt{assignable} member when its only entry is \texttt{\textbackslash nothing}
and interprets an absent member as \texttt{\textbackslash nothing} on read,
saving one array entry per specification. A specification with genuine frame
clauses lists them explicitly.

The third is a \emph{token dictionary}. A twelve-entry dictionary substitutes
the most common JML tokens --- \texttt{\textbackslash result},
\texttt{\textbackslash old(}, the quantifiers, the implication arrow, the
null comparisons, the frame keywords --- with single Unicode control bytes that
cannot appear in legitimate JML, so the substitution is unambiguous and
invertible. Because UTF-8 encodes each control byte as a single byte, each token
occurrence saves the token's length minus one. The reader decodes on read; a
consumer ignorant of the dictionary sees the control characters and can either
decode using the published table or fall back to the sidecar.

The choice of which tokens to encode and which clause to default is not
arbitrary; each rests on the shape distribution measured on the real-inference
corpus (Section~\ref{sec:dist-results}). A format designed without the
measurement would either include every JML keyword in the dictionary, paying
constant-pool cost for tokens that never fire, or omit tokens that occur
frequently in real specifications. The divergence between the optimisation's
effect on synthetic and on real specifications, reported below, is itself
evidence that the saving depends on which clause shapes the consumer's inferrer
produces.

\subsection{Canonical Form}

To support byte-level roundtrip equality, clause text is canonicalised before
emission. Whitespace runs collapse to single spaces; there is no space around
member-access or method-reference operators; binary operators take a single
surrounding space; there is no space inside brackets; and the short implication
and equivalence arrows are normalised to their long forms. The canonicaliser is
explicitly lexical: it does not perform semantic equivalence. Two clauses that
differ in operand order remain distinct strings, and roundtrip equality is
preserved on the wire form rather than on a normalised theory level. This is a
deliberate choice. Semantic equivalence would silently rewrite valid
specifications in transit, which is unacceptable when the embedder's job is to
carry data faithfully; a consumer that wants semantic normalisation can apply it
after recovery, but the transport layer must not.

\subsection{The OpenAPI Extension and the Sidecar}

For service boundaries, a parallel extension family carries the same content in
an OpenAPI~3.x document: \texttt{x-jml-requires}, \texttt{x-jml-ensures},
\texttt{x-jml-assignable}, \texttt{x-jml-signals}, and \texttt{x-jml-invariant},
each an array of canonical-form strings, carried per operation and per schema.
The \texttt{x-} prefix is the OpenAPI mechanism for vendor extensions, which
standard parsers preserve without understanding; the extension therefore rides
through the existing OpenAPI toolchain unmodified. A JSON Schema document
accompanies the artefact for independent validation, and where the OpenAPI
document is generated by a framework and not editable, a sidecar JSON file
carries the same content keyed by route and verb.

For specifications that exceed the per-class bytecode budget --- rare, but
possible for a method with many quantified loop invariants --- and for signed
JARs that cannot be re-signed, the embedder emits a parallel artefact carrying
JML stub files keyed by class internal name, with a manifest entry recording the
source JAR's hash so consumers can detect drift. The stub files are the format
OpenJML reads natively, so the sidecar is consumable by the verifier without
further transformation. The sidecar is the universal fallback: where in-bytecode
embedding is unsuitable, the same content travels alongside.

\subsection{The OpenAPI Extension in Depth}

The service-side extension warrants a fuller treatment, because a service boundary
differs from a library boundary in ways that shape the design. A library's contract
attaches to a method, identified by its name and bytecode descriptor; a service's
contract attaches to an \emph{operation}, identified by a route and an HTTP verb,
and to the \emph{schemas} of the data it exchanges. The extension therefore carries
contracts at two levels: per-operation, where \texttt{x-jml-requires} states the
preconditions on a request and \texttt{x-jml-ensures} the postconditions on a
response, and per-schema, where \texttt{x-jml-invariant} states the invariants a
data type maintains. The two levels mirror the method-and-class structure of the
bytecode side, adapted to the operation-and-schema structure of a service
interface.

The \texttt{x-jml-signals} key handles exceptional behaviour, which at a service
boundary means error responses. It carries an array of exception-condition pairs,
the exception named by its JVM internal name so that a cross-language consumer can
resolve it to whatever its own language calls the corresponding error. This
cross-language resolvability matters because a service's clients are often written
in languages other than the service's, and a contract that named exceptions in the
service's language alone would be opaque to them; the internal name is a
language-neutral identifier the client can map.

The extension's compatibility rests on the OpenAPI specification's own provision
for vendor extensions: any key prefixed \texttt{x-} is reserved for extension data,
and conforming parsers must preserve it without interpreting it. The extension
therefore rides through the standard OpenAPI toolchain --- the parsers, the client
generators, the documentation renderers --- untouched, recovered by a consumer that
knows the \texttt{x-jml-*} family and ignored by one that does not. A JSON Schema
document accompanies the extension so that a consumer can validate an OpenAPI
document's contracts independently of the inferrer, and where the OpenAPI document
is generated by a framework and not hand-editable, the same content travels in a
sidecar JSON file keyed by route and verb, the service-side analogue of the bytecode
sidecar. The design thus extends the toolchain-transparency property from the
bytecode boundary to the service boundary by the same strategy: ride the format's
own extension mechanism, and fall back to a sidecar where the artefact cannot be
modified.

\section{Deployment Scenarios}
\label{sec:dist-scenarios}

The format design admits three deployment paths, and choosing among them is a
matter of matching the carrier to the constraints of the artefact. The default
path is in-bytecode embedding: the inferrer runs over a library's source, the
embedder writes the specifications into the compiled JAR's class files, and the
enriched JAR is published in place of the original. This is the path for the
common case --- an unsigned library whose specifications fit the per-class budget
--- and it is the path the empirical validation of Section~\ref{sec:dist-results}
exercises. The consumer receives an ordinary-looking JAR that happens to carry
recoverable contracts, and tooling that knows to look for them finds them while
tooling that does not is unaffected.

The second path is the OpenAPI extension, for service boundaries. A service that
documents its interface with an OpenAPI document --- the dominant convention for
REST APIs --- can carry per-operation and per-schema contracts in the
\texttt{x-jml-*} extension family, so that a client generator or a contract-test
tool reading the document recovers the same preconditions and postconditions a
library consumer would recover from bytecode. The path matters because the
service boundary is where the oracle problem is most acute: a client of a remote
service has neither its source nor its bytecode, only its interface description,
and a contract carried in that description is the only specification the client
can obtain. The OpenAPI path extends the reach of inferred specifications from
the library-integration scenario of Chapter~\ref{ch:introduction} to the
service-integration scenario, and it is the mechanism the future-work discussion
of service-boundary testing relies on.

The third path is the sidecar JAR, the universal fallback. Two constraints force
it. A signed JAR cannot be modified without invalidating its signature, and
re-signing requires the signer's key, which the party enriching the JAR may not
hold; embedding into such a JAR is impossible, but a parallel sidecar artefact
carrying the same specifications as JML stub files travels alongside it without
touching the signed bytes. And a specification too large for the per-class
bytecode budget --- rare, but possible for a method bearing many quantified loop
invariants --- can overflow the annotation, in which case the sidecar carries the
overflow. The sidecar's stub files are the format OpenJML reads natively, so the
fallback is directly consumable by the verifier, and the manifest records the
source JAR's hash so a consumer can detect drift between the sidecar and the
artefact it describes. The three paths together ensure that no artefact is
un-enrichable: where in-bytecode embedding is unsuitable, the same content
travels as a sidecar, and where the artefact is a service rather than a library,
the same content travels in the interface description.

\section{Implementation}
\label{sec:dist-impl}

The embedder is built on the ASM bytecode-engineering library, which permits
class-file attributes to be added to an existing JAR in place, without
recompilation. The writer walks each class with an ASM \texttt{ClassVisitor},
and for each method, constructor, field, and type carrying a specification, it
emits the packed annotation into the element's
\texttt{RuntimeVisibleAnnotations} attribute. The reader is the inverse: it
walks the embedded class files and recovers the specification from the
annotation, decoding the token dictionary and synthesising the omitted frame
default.

A subtlety in the ASM visitor lifecycle proved to be the source of the most
consequential bug in the implementation, and it is instructive. The writer's
lazy annotation emission was originally triggered by the visitor callbacks that
fire for methods with bodies --- the callbacks for code, parameters, and
existing annotations. None of these fires for abstract or interface methods,
which have no body and invoke only the end-of-method callback. The result was
that interface methods silently received no embedded annotation, and an early
prototype reported a roundtrip rate of 96.9\% on one library and 97.6\% on
another, the shortfall concentrated entirely in interfaces. Adding the
end-of-method callback to the trigger set closed the gap to 100\%. The episode
is a reminder that the cases a mechanism handles silently --- here, methods
without bodies --- are exactly the cases where a measurement that does not seek
them out will miss the failure.

The bug is worth dwelling on because of what it reveals about measurement design.
The roundtrip rate of 96.9\% looked like success --- a small residual on a hard
problem --- and a study satisfied with a high-but-imperfect number would have
reported it and moved on, attributing the gap to some intractable corner. The gap
was not intractable; it was a single missing callback, concentrated in a
structural feature of the code (interfaces) that the aggregate rate obscured.
Disaggregating the roundtrip failures by program element revealed that every
failure was an interface method, which pointed directly at the lifecycle callback
that interfaces alone exercise. The lesson the thesis draws is that an imperfect
aggregate should be disaggregated before it is accepted, because a residual
concentrated in one structural category is usually a bug with a specific cause
rather than an inherent limit. The 100\% rate that the disaggregation made possible
is a stronger result than the 96.9\% would have been, and it was available only
because the measurement sought out where the failures fell rather than tolerating
their rate.

Multi-release JARs --- archives carrying version-specific class files for different
JDK versions in parallel directory trees --- are the other structural feature that
stressed the embedder, because a naive walk would embed into only one version's
classes or would mismatch versions. The embedder handles the multi-release layout
by embedding into each version's class files independently, so a consumer running
under any supported JDK recovers the specification from the class file that JDK
loads. Guava, a multi-release JAR, exercises this path, and its 100\% roundtrip
rate confirms the handling; the threats section notes that ahead-of-time
compilation, which transforms the class files before they reach a JVM, is the
remaining packaging variant the evaluation does not cover.

\section{What Lossless Roundtrip Means, Precisely}
\label{sec:dist-lossless}

The central correctness property of the embedder is \emph{lossless roundtrip}:
the specification recovered by the reader equals the specification the writer was
given. Stating the property precisely matters, because there is a weaker property
the embedder deliberately does not claim and a stronger one it deliberately does
not need. The embedder claims that for every method, the clause list the reader
recovers is equal, as a list of canonical-form strings, to the clause list the
writer embedded. It does not claim \emph{semantic} equality --- that the recovered
clauses are logically equivalent to the originals --- because semantic
equivalence is both unnecessary for faithful transport and dangerous to attempt.

The danger is worth dwelling on. A roundtrip that normalised semantics --- that
recognised \texttt{a > b} and \texttt{b < a} as equal and stored a canonical
representative --- would silently rewrite a consumer's specification in transit.
A consumer that wrote \texttt{a > b} and recovered \texttt{b < a} would have its
data altered, however harmlessly, by a layer whose only job is to carry it
faithfully. Worse, semantic normalisation is undecidable in general and unsound in
the corners where it is decidable, so an embedder attempting it would introduce a
class of bugs --- specifications subtly changed in transit --- that a purely
lexical embedder cannot have. The embedder therefore canonicalises only
\emph{lexically} (Section~\ref{sec:dist-format}): it fixes whitespace, operator
spacing, and arrow forms, producing a single well-defined wire string for any
given clause, but it never rewrites the logical content. Two logically equivalent
clauses that differ lexically remain two distinct strings, and the roundtrip
preserves each exactly.

This lexical discipline is what makes the 100\% roundtrip claim both achievable
and meaningful. Achievable, because lexical equality is decidable and the
canonical form makes it stable across the write-read cycle; meaningful, because a
consumer receives exactly the bytes the producer's inferrer emitted, with no
silent transformation. The property is verified empirically by embedding and
re-reading every specification and comparing clause lists for exact string
equality, which is the measurement Section~\ref{sec:dist-results} reports. The
distinction between this lexical roundtrip and a semantic one is not pedantry: it
is the difference between a transport layer that can be trusted to carry data
unchanged and one that cannot, and the choice of the lexical property is what
earns the embedder the right to be called lossless.

\section{Empirical Validation}
\label{sec:dist-results}

The embedder is evaluated on real-world Java libraries under two regimes. The
first uses \emph{synthetic} specifications --- one precondition per non-primitive
parameter, one postcondition for reference returns, and the default frame ---
to isolate the embedding mechanism from the inferrer's behaviour and stress it
uniformly. The second uses \emph{real inferred} specifications produced by
running JML-Inferrer over the libraries' sources, exercising the full
inference-to-embedding pipeline including quantifiers, history references, and
branch-conditional implications. The headline results appear in
Table~\ref{tab:dist-validation}.

\begin{table}[ht]
\centering
\caption{Embedder validation on real-world libraries. Method counts,
lossless-roundtrip rate, and overhead are deterministic and reported once;
throughput is the median across 39 runs. The fourth-row group reports the full
inferrer-to-embedder pipeline on real inferred specifications.}
\label{tab:dist-validation}
\begin{tabular}{lrrrrr}
\toprule
\textbf{Library} & \textbf{Methods} & \textbf{Lossless} & \textbf{Byte ovhd.} & \textbf{Embed (m/s)} & \textbf{Read (m/s)} \\
\midrule
\multicolumn{6}{l}{\textit{Synthetic specs (isolating the embedding mechanism)}} \\
Apache Commons Lang 3.14.0 & 4{,}029 & 100.0\% & +5.85\% & 16{,}225 & 221{,}617 \\
Apache Commons IO 2.13.0 & 2{,}677 & 100.0\% & +6.97\% & 21{,}976 & 327{,}618 \\
Guava 33.3.0-jre & 13{,}840 & 100.0\% & +5.27\% & 12{,}618 & 127{,}924 \\
\midrule
\multicolumn{6}{l}{\textit{Real inferred specs (full inferrer$\to$embedder pipeline)}} \\
Apache Commons Lang 3.14.0 & 2{,}332 & 100.0\% & +10.07\% & 8{,}648 & 103{,}779 \\
Apache Commons IO 2.13.0   & 1{,}680 & 100.0\% & +10.20\% & 10{,}215 & 107{,}318 \\
Guava 33.3.0-jre           & 6{,}320 & 100.0\% & +4.63\%  & 6{,}052 & 74{,}144 \\
JML-Inferrer self-test     & 73    & 100.0\% & --- & --- & --- \\
\bottomrule
\end{tabular}
\end{table}

\subsection{Lossless Roundtrip (RQ3a)}

Every corpus achieves a 100.0\% lossless-roundtrip rate: every specification the
harness embeds is recovered with byte-for-byte equality on its clause list, with
zero mismatches anywhere. The real-inference regime is the more demanding test,
because it exercises every JML construct the engine produces rather than the
synthetic worst case. Across the three libraries the real-inference harness
processes 1{,}118 source files containing 16{,}507 source method declarations,
resolves 10{,}332 of them to distinct bytecode keys, and roundtrips every one
losslessly. The constructs exercised include nested quantifiers, history
references via \texttt{\textbackslash old}, conditional postcondition
implications, and loop invariants --- the structurally complex output, not
merely the simple null-check preconditions of the synthetic case. The
distribution of constructs is reported in Table~\ref{tab:dist-shapes}; every
category roundtrips losslessly on every library.

\begin{table}[ht]
\centering
\caption{Spec-shape coverage on the 2{,}332 Commons Lang specifications embedded
through the real-inference harness (representative; the other libraries give
similar distributions). Every category roundtrips losslessly.}
\label{tab:dist-shapes}
\begin{tabular}{lrr}
\toprule
\textbf{Construct} & \textbf{Specs} & \textbf{\% of corpus} \\
\midrule
\texttt{requires} clauses               & 1{,}116 & 47.9\% \\
\texttt{ensures} clauses                & 2{,}109 & 90.4\% \\
\texttt{assignable} clauses             & 2{,}332 & 100.0\% \\
\texttt{loop\_invariant} clauses        & 269 & 11.5\% \\
Quantifiers (\texttt{\textbackslash forall}/\texttt{\textbackslash exists}/\texttt{\textbackslash sum}) & 16 & 0.7\% \\
\texttt{\textbackslash old} references  & 85 & 3.6\% \\
\texttt{\textbackslash result} references & 1{,}725 & 74.0\% \\
Branch-conditional (\texttt{==>})       & 345 & 14.8\% \\
\bottomrule
\end{tabular}
\end{table}

\subsection{Byte Overhead (RQ3b)}

Synthetic-spec overhead sits between 5.27\% and 6.97\%; the differences across
libraries are explained by structural shape rather than total size, with Commons
IO running highest because its higher proportion of interfaces means more methods
each carrying a clause set. In absolute terms the synthetic overhead is 38~KB
for Commons Lang, 34~KB for Commons IO, and 162~KB for Guava --- well below
typical distribution sizes.

The synthetic figures are not an upper bound on real overhead, a point worth
emphasising because a deployer planning a storage budget needs the realistic
figure. The real-inference run on the same Commons Lang JAR measures 10.07\%,
more than four percentage points above the synthetic 5.85\%, because real
inferred specifications are richer than the synthetic worst case: 90.4\% carry an
\texttt{ensures} clause against the synthetic one-per-reference-return policy,
11.5\% carry loop invariants the synthetic harness never emits, 14.8\% include
branch-conditional implications, and the average clause string is longer because
the engine emits parameter bounds, result relations, and history references the
synthetic policy never produces. Guava inverts the relationship --- its
real-inference overhead, 4.63\%, is \emph{lower} than its synthetic figure ---
because its large generic-heavy classes have dense bodies, so the same per-spec
byte cost amortises against a larger denominator. The lesson for a deployer is to
provision for the real-inference figure, around ten percent on a typical utility
library, not the synthetic-uniform figure.

The Guava inversion is worth dwelling on because it corrects an intuition that
would otherwise mislead a deployer. The natural expectation is that real
specifications, being richer than the synthetic worst case, always cost more, and
on the smaller utility libraries they do --- four to five points more. Guava
violates the expectation not because its specifications are leaner but because its
denominator is larger: overhead is the embedded bytes divided by the original
bytes, and Guava's original bytes are inflated by large generic-heavy classes with
dense method bodies, so the same per-method specification cost is a smaller
fraction of a larger whole. The general principle the inversion reveals is that
overhead depends on two quantities a deployer should consider separately --- the
absolute specification cost per method, which is a property of the inferrer's
output, and the original code size per method, which is a property of the library's
density --- and that a dense library can carry richer specifications at lower
proportional overhead than a sparse one. The synthetic-versus-real comparison
isolates the first quantity by holding the specification policy fixed, and the
cross-library comparison isolates the second by holding the policy fixed and
varying the library, and the Guava inversion is what happens when the second effect
dominates the first.

\subsection{Format Optimisation}

Table~\ref{tab:dist-opt} reports the impact of the three size optimisations
against the verbose predecessor format. The optimisations reduce overhead by 48
to 55 percent on the synthetic corpora and by 31 percent on real inference,
while preserving the 100\% lossless property in every case.

\begin{table}[ht]
\centering
\caption{Format-optimisation impact on bytecode-size overhead. The
lossless-roundtrip rate is 100\% under both formats.}
\label{tab:dist-opt}
\begin{tabular}{lrrr}
\toprule
\textbf{Corpus} & \textbf{Verbose ovhd.} & \textbf{Optimised ovhd.} & \textbf{Reduction} \\
\midrule
Synthetic --- Commons Lang      & 11.95\% & 5.85\%  & $-$51\% \\
Synthetic --- Commons IO        & 13.49\% & 6.97\%  & $-$48\% \\
Synthetic --- Guava             & 11.76\% & 5.27\%  & $-$55\% \\
Real inference --- Commons Lang & 14.70\% & 10.07\% & $-$31\% \\
\bottomrule
\end{tabular}
\end{table}

The synthetic reduction is larger than the real-inference reduction because
synthetic specifications are dominated by exactly the patterns the optimisation
targets --- every specification carries the default-omittable frame, every
precondition is a null check, every postcondition references the result --- all
of which the optimisation captures. Real specifications contain a wider variety
of clause shapes, many of which the twelve-token dictionary does not cover, so
the real-inference reduction is the more realistic estimate for a production
deployment. A configuration sweep confirms the optimisations' contribution: an
always-explicit mode with no codec runs 1.19 percentage points above the
default, and an automatically learned twenty-eight-token dictionary gains a
further 0.59 points over the hand-picked twelve, establishing that the
hand-picked dictionary already captures most of the available saving. A per-spec
DEFLATE compression mode is a clear negative result, running 20 points
\emph{worse} than the default, because DEFLATE cannot exploit redundancy in the
fifty-to-hundred-byte payloads typical after token encoding and the per-spec
compressed payloads forfeit the constant-pool deduplication the uncompressed
format enjoys. The negative result is reported because it is informative: it
shows that the right granularity for compression, if any, is the class file
rather than the individual specification.

\subsection{Throughput (RQ3c)}

Embedding throughput ranges across runs from roughly six thousand to thirty-one
thousand methods per second, and reading throughput an order of magnitude higher,
from roughly seventy-four thousand to over three hundred thousand. At the
medians, embedding a fourteen-thousand-method JAR completes in well under a
second end to end, and reading completes faster still because the reader skips
full class rewriting. These figures comfortably support per-pull-request
continuous-integration workflows: even a pessimistic re-run on every push of a
Guava-sized library is sub-second, far below the Maven build time the inference
pipeline already rides on. The throughput question is therefore answered
affirmatively with room to spare.

\subsection{Throughput in Context (RQ3c)}

The throughput question is answered affirmatively, but the figures repay
interpretation against the workflows they are meant to support. Embedding a
fourteen-thousand-method library completes in well under a second at the median
embedding rate, and reading completes an order of magnitude faster because the
reader does not rewrite the class files. The relevant comparison is not against
some absolute standard but against the cost of the surrounding workflow. The
embedding rides on a Maven build that already takes tens of seconds to minutes;
the inference that produces the specifications takes longer still, at tens of
milliseconds per file over thousands of files. Against this backdrop the sub-second
embedding is negligible --- it disappears into the noise of the build it
augments.

This matters for the continuous-integration scenario the question targets. A
plausible deployment runs the inferrer and the embedder on every pull request, so
that the enriched artefact reflects the current source. For this to be acceptable
the embedding must fit within the per-pull-request budget, which is dominated by
compilation and testing and measured in minutes. A sub-second embedding step fits
trivially, even re-run from scratch on every push of the largest library measured.
The reading throughput matters for the consumer side: a tool that recovers
specifications from a dependency does so at hundreds of thousands of methods per
second, fast enough that recovering an entire library's contracts is imperceptible.
The throughput figures thus clear their bar with orders of magnitude to spare,
which is the right outcome to report --- not a tight fit that might fail under
load, but a comfortable one that will not.

The variation in throughput across runs --- the reason the figures are reported as
medians over many runs rather than as single measurements --- reflects the ordinary
sources of variance in JVM workloads: just-in-time compilation warming up, garbage
collection pausing, the operating system scheduling other work. These sources
affect the absolute rate but not the conclusion, because even the slowest observed
run completes the embedding within the per-pull-request budget. Reporting the
median over many runs, with the observed range, is the honest way to characterise a
throughput that varies for reasons unrelated to the embedder's logic, and it is the
same discipline the lossless-roundtrip and overhead measurements do not need
because those are deterministic.

\subsection{Toolchain Compatibility (RQ3d)}

The compatibility question is answered on two fronts. On the bytecode side, a
negative test loads an embedded class through a class loader that does not see
the annotation type, and confirms that the class loads successfully:
reflective annotation access either returns the resolvable annotations or raises
the documented exception, but class loading itself is unaffected. A JVM with no
awareness of the annotation type runs the embedded class exactly as it would the
unembedded one --- the embedding is invisible to the runtime. On the service
side, the OpenAPI extension is preserved by the standard parsers without
modification, because the \texttt{x-} prefix is the specification's own
mechanism for vendor extensions. The mechanism therefore meets its design goal:
enrichment requires no change to how artefacts are built, loaded, published, or
parsed.

\section{The Class-File Format and Why the Mechanism Is Transparent}
\label{sec:dist-classfile}

The transparency of the embedding rests on a feature of the Java class-file
format that is worth making explicit, because it is what allows enrichment
without toolchain modification. A Java class file is a structured binary
container: a magic number, version, constant pool, and a sequence of structures
describing the class's fields, methods, and attributes. The format is
deliberately extensible through \emph{attributes} --- named, length-prefixed
blobs attached to the class, its fields, and its methods --- and the
specification mandates that a reader encountering an attribute it does not
recognise must skip it, using the length prefix, rather than fail. This
skip-the-unknown rule is the foundation of forward compatibility in the JVM, and
it is what guarantees that a class file carrying an unfamiliar attribute loads
correctly in a JVM that has never heard of it.

Runtime-visible annotations are carried in exactly such an attribute, the
\texttt{RuntimeVisibleAnnotations} attribute, which the format defines and every
conforming JVM reads. An annotation written into this attribute is recoverable
by reflection if the annotation's type is on the classpath, and is harmlessly
skipped if it is not. The embedder exploits this: it writes the packed
specification annotation into the \texttt{RuntimeVisibleAnnotations} attribute of
each method, field, and type. A consumer that has the annotation type can recover
the specification by reflection; a class-file reader that lacks it (such as the
ASM-based reader the extractor uses) can recover it by reading the attribute
directly; and a JVM that knows nothing of it loads and runs the class exactly as
before. The negative-load test of Section~\ref{sec:dist-results} verifies the
last case empirically, and the class-file format's skip-the-unknown rule is the
reason it passes.

The constant pool is the other relevant feature. Strings in a class file are
stored once in the constant pool and referenced by index, so two methods whose
specifications share a clause string --- and inferred specifications share clause
strings heavily, since \texttt{requires arr != null} and \texttt{assignable
\textbackslash nothing} recur across a library --- pay for that string only once.
This deduplication is what makes the per-method annotation overhead modest in
aggregate, and it is precisely the saving that the per-spec DEFLATE mode of
Section~\ref{sec:dist-results} forfeits by making each method's payload unique,
which is the mechanism behind that mode's negative result. The format's existing
machinery for sharing strings does, for free, what per-spec compression cannot.

\subsection{A Worked Embedding}

Consider the \texttt{max} method of Chapter~\ref{ch:background}, with its
inferred contract of one trivial precondition, two postconditions, and the
default frame. The embedder writes a single \texttt{@JmlSpecs} annotation onto
the method, with \texttt{ensures} carrying the two postcondition strings and the
other members defaulting --- \texttt{requires true} is the trivial precondition
the format omits, and \texttt{assignable \textbackslash nothing} is the default
the format synthesises on read. The two postcondition strings are token-encoded:
the \texttt{\textbackslash result} occurrences collapse to a single control byte
each, and the strings are interned in the constant pool, so the second
postcondition's shared substructure with the first costs nothing beyond its
distinct portion. On read, the extractor recovers the annotation, decodes the
token bytes back to \texttt{\textbackslash result}, synthesises the omitted frame
and trivial precondition, and reconstructs the contract byte-for-byte. The
roundtrip is lossless because the canonical form fixed the wire representation
before emission, so there is a single, well-defined string for the reader to
reproduce. This worked case is the unit that the 100\% lossless-roundtrip rate
of Section~\ref{sec:dist-results} aggregates over twenty thousand times.

\section{Consuming Embedded Specifications}
\label{sec:dist-consume}

The embedding is only useful if consumers can act on it, and the format is
designed so that several classes of consumer can, each by a mechanism natural to
it. A reflective consumer --- a test generator or a runtime contract checker
running inside a JVM that has the annotation type on its classpath --- recovers
the specification by ordinary reflection, reading the \texttt{@JmlSpecs}
annotation off the method and decoding its members. This is the path for a tool
that wants to enforce contracts at runtime or to drive test generation from them,
and it requires nothing of the consumer beyond depending on the small annotation-
type library.

A static consumer --- a verifier, an IDE, an analysis tool --- recovers the
specification by reading the class file directly, without loading it into a JVM.
The ASM-based reader the extractor uses is one such consumer, and any class-file
reader can be taught the format from the specification of
Section~\ref{sec:dist-format}. This is the path for OpenJML itself: a verifier
checking client code against a dependency can read the dependency's embedded
contracts and use them as the assumed specifications of the called methods,
exactly as it would use hand-written stubs, closing the modular-verification loop
across the library boundary without the library's source.

A language-model consumer --- the case that motivates much of the thesis ---
recovers the specification however its tooling extracts class-file metadata, and
incorporates it into the model's context as the inferred contract of the API under
discussion. This is the mechanism by which the test-generation benefit of
Chapter~\ref{ch:testgen} reaches a model working against a compiled dependency:
the model reads the embedded contract rather than guessing the dependency's
behaviour, and the same improvement the controlled experiment measured on
source-available code becomes available on compiled code. The service-side OpenAPI
extension serves the same role for a model reasoning about a remote service from
its interface description.

The robustness of the consumer story rests on the property the negative-load test
establishes: an embedded artefact is safe to ship to consumers who know nothing of
the embedding. Were the embedding to break class loading in a JVM without the
annotation type, every consumer would have to add the annotation library to its
classpath before using the enriched artefact at all --- precisely the toolchain
coupling the design exists to avoid. The test confirms the coupling does not arise:
the class-file format's skip-the-unknown rule tolerates a runtime-visible
annotation whose type cannot be resolved, so the worst that happens is that
reflective access to the annotation raises the documented type-not-present
exception, which a consumer that does not ask for the annotation never triggers.
The class loads, the methods run, and a consumer unaware of the embedding is
unaffected --- the precondition for enriching an artefact that will be consumed by
parties the enricher does not control. The JSON Schema document accompanying the
OpenAPI extension serves the analogous role on the service side, letting a consumer
validate that the contracts it receives are well-formed before acting on them, the
way it would validate any external input.

The forward-compatibility of the format matters for all three consumer classes. A
consumer built against the current clause kinds skips members it does not
recognise, so a future clause kind added to the format does not break existing
consumers; and a consumer that does not know the token dictionary sees the control
bytes and can either decode them from the published table or fall back to the
sidecar's uncompressed stub files. The format degrades gracefully: a consumer
recovers as much of the specification as it understands and ignores the rest, which
is the property that lets the format evolve without coordinating every consumer's
upgrade.

\section{Discussion and Threats}
\label{sec:dist-discussion}

\subsection{What the Result Establishes}

The contribution is a path from an inferred specification to a consumer-readable
contract that requires no change to the build toolchain: any existing JAR can be
enriched in place, and the enrichment survives ordinary distribution and is
recoverable by any class-file reader without the source. The 100\% lossless rate
across four corpora and every JML construct the engine produces establishes that
the transport is faithful; the sub-eleven-percent overhead and sub-second
throughput establish that it is affordable; the negative-load test and the
OpenAPI parser compatibility establish that it is transparent. Together these
close the distribution gap that the test-generation result of
Chapter~\ref{ch:testgen} would otherwise leave open: a specification that
improves downstream tooling is now available to downstream consumers who lack
the source, which is most of them.

The result also completes the integration scenario of
Chapter~\ref{ch:introduction}. A developer adopting a compiled library, or a model
reasoning about one, encounters it as bytecode without a contract, and learns its
behaviour by trial, by reading whatever source is available, or by the exceptions a
violated precondition throws. With the embedding, the library's inferred contract
travels inside the JAR, and tooling at the call site --- an IDE, a verifier, a
contract-checking test generator, a model reading the API --- recovers it without
the source. The integration scenario's central difficulty, that the contract is
present in the library's design but absent from the form the integrator encounters,
is resolved by carrying the contract into that form. The distribution study is, in
this sense, the bridge between the inference study, which produces a contract from
source, and the consumer scenarios, which need the contract without source: it is
the stage that makes the produced contract reach the consumer that the value study
showed would benefit from it.

\subsection{Format Evolution and Versioning}

A distribution format must evolve, and the design anticipates evolution in a way
that the empirical results validate. The format carries a version member, and the
reader accepts both the current packed format and its verbose predecessor, so that
artefacts embedded against the earlier format continue to load against the current
reader. This backward compatibility is not merely good practice; it is a
correctness property the thesis tests, because a format that broke its predecessor
on every revision would force every consumer to upgrade in lockstep, which is
exactly the toolchain coupling the design aims to avoid.

The forward direction is handled by the skip-the-unknown discipline of the
class-file format (Section~\ref{sec:dist-classfile}) and the skip-unrecognised-
members discipline of the annotation. A future clause kind, added as a new
annotation member, is ignored by a reader that predates it rather than breaking
that reader; a future format version is distinguished by its version member, and a
reader can dispatch on it. The format can therefore add clause kinds and revise
its encoding without invalidating either the artefacts embedded against earlier
versions or the consumers built against them, which is the property that lets the
inference engine's output grow --- new clause kinds as new patterns are added ---
without a coordinated upgrade of every consumer in the ecosystem.

The configuration sweep of Section~\ref{sec:dist-results} bears on evolution in a
second way. By exposing the default-omission and token-dictionary conventions as
configurable, the format lets a deployer whose inferrer has a different dominant
clause shape opt out of conventions tuned for this thesis's engine and supply its
own. An inferrer that emitted a conservative \texttt{assignable \textbackslash
everything} default rather than \texttt{\textbackslash nothing} would select a
different default-omission policy; an inferrer with a different token frequency
profile would supply a different dictionary, or learn one from its own corpus with
the dictionary learner the sweep evaluates. The format is thus not hard-coded to
one inferrer's output distribution but parameterised by it, which is what allows
the same transport mechanism to serve inferrers other than the one this thesis
built.

\subsection{When to Embed and When Not To}

The mechanism's availability does not make embedding always the right choice, and a
deployer should weigh the considerations the evaluation surfaces. Embedding is the
right default when the consumer benefits from structured contracts and the artefact
can be modified --- an unsigned library published to a repository, enriched once at
build time and consumed by tooling that recovers the contracts. The overhead, under
eleven percent on a typical library, is negligible against distribution sizes, and
the throughput fits any continuous-integration budget, so the cost side of the
ledger is light.

Embedding is the wrong choice in three situations the design anticipates. A signed
artefact that cannot be re-signed forecloses in-bytecode embedding, and the sidecar
is the only option. An artefact whose consumers are uniformly unaware of
specifications gains nothing from embedding and pays the overhead for nothing,
though the overhead is small enough that the cost is more theoretical than real. And
an artefact where even the small overhead is unacceptable --- a tightly
size-constrained deployment --- should use the sidecar, which moves the bytes out of
the artefact entirely. The configuration sweep gives a deployer the further choice of
trading overhead against codec complexity: the always-explicit mode pays a little
more for the simplicity of no token dictionary, and the learned-dictionary mode pays
a little less at the cost of carrying a corpus-specific dictionary.

The decision, in short, is not whether the mechanism works --- the evaluation
establishes that it does, losslessly and cheaply --- but whether the consumer will
use the contracts and whether the artefact can carry them. The thesis provides the
mechanism and the measurements; the deployment decision rests on the consumer
landscape and the artefact constraints, which vary by context. The design's
provision of three paths --- in-bytecode, OpenAPI, sidecar --- is what ensures the
decision is always ``which path'' rather than ``whether possible''.

\subsection{Threats to Validity}

The evaluation covers three libraries plus a self-test rather than the full
protocol of ten; the lossless analysis identifies the residual risks ---
interface methods and multi-release JARs --- and the embedder reaches 100\% on
libraries exercising both, so the remaining libraries are mechanical to validate.
The evaluation targets JDK~21; older JDKs and ahead-of-time compilers such as
GraalVM native-image are deferred, though the class files emitted by Commons Lang
target Java~8, so the runtime-side validation should generalise. The source-to-
bytecode key resolution in the real-inference harness matches on name and arity
only, so a residual of methods is lost to overload collisions where two source
overloads map to the same erased descriptor; the 100\% lossless rate is on every
specification the harness successfully embeds, so the residual is a measurement
limitation rather than an embedder limitation, and a full type-resolving harness
would recover it without changing the embedder's behaviour. These threats bound
the generality of the measurement but not the central claim, which is that the
embedding mechanism transports every specification it is given, losslessly, at
acceptable cost, through an unmodified toolchain.

A reflection on the nature of this study closes the chapter. Unlike the
test-generation study, whose claim is causal and whose validation is statistical,
the distribution study's claims are about a deterministic mechanism, and its
validation is exact measurement rather than inference. There is no sampling
uncertainty in a losslessness rate: a specification either roundtrips or it does
not, and the rate is a count, not an estimate. This gives the distribution study a
different epistemic character --- its claims are not ``probably true within a
confidence interval'' but ``observed to be true on every case measured'' --- and the
threats that bound it are about coverage (which cases were measured) rather than
about uncertainty (how confident the estimate is). The two studies' different
characters reflect the different kinds of claim they make, and matching the
validation to the claim's kind, as Chapter~\ref{ch:introduction} argued, is what
lets each study assert exactly the certainty its claim warrants: the value study a
confident causal effect, the distribution study an exact mechanical property, and
neither overreaching into the other's kind of assertion.
